%% kapa: exam=1 round=25 year=2014 period=1 form=A subject=economics subject_ko=경제원론 seq=12 number=52
%% source: Q-Net 감정평가사 기출문제 (cst00309) artlSeq=5101207 file=extracted/1008/5101207/2014년 제25회 감정평가사 1차시험 기출문제/1교시(A형).hwp
%% conversion: kordoc 4.24.51 -> tools/gen1.py (auto); review=proofread-sonnet; flags=none
\begin{question}{52}
효용을 극대화하는 근로자 甲은 여가와 근로소득을 선택한다. 다음 중 관찰될 수 있는 경우를 모두 고른 것은? (단, 甲에게 여가는 정상재이다.)
\begin{framedtext}
ㄱ. 시간당 임금이 상승했는데, 甲의 노동공급이 감소했다. \par
ㄴ. 시간당 임금이 상승했는데, 甲의 노동공급이 증가했다. \par
ㄷ. 시간당 임금에 근로소득세를 부과했더니, 甲의 노동공급이 증가했다. \par
ㄹ. 甲에게 비근로소득이 생겨 노동공급이 증가했다. \par
\end{framedtext}
\begin{choices}
  \item ㄱ, ㄴ
  \item ㄴ, ㄹ
  \item ㄱ, ㄴ, ㄷ
  \item ㄱ, ㄷ, ㄹ
  \item ㄴ, ㄷ, ㄹ
\end{choices}
\end{question}
\answer{3}
